\section{Introduction}

Machine learning classifiers can change their prediction under small input perturbations, raising concerns about safety, security, and ethics. 
A central notion for studying this phenomenon is \emph{global robustness}, which requires that similar inputs receive identical (or similar) outputs~\cite{casadio2022neural}. 
Similarity is captured by a distance function, making the framework remarkably general: in computer vision it led to weaker notions, since similarity based on pixel intensities propagates by transitivity; but
in fairness, where the distance encodes context-dependent similarity of individuals under ethical constraints (e.g., equally qualified for a job)~\cite{fairness-landscape,kusner2017counterfactual,ranking-fairness}, and in natural language processing, syntactic and semantic similarity is encoded under meaningful linguistic variation~\cite{levy2025statistical} remains central.
Importantly, the notion of similarity may \emph{not} be metric, violating the triangle inequality~\cite{garcia2019learning,similarity-guide}.
Motivated by this observation, we study robustness over arbitrary \emph{semimetric} spaces.
A common observation is that exact global $\epsilon$-robustness is essentially unsatisfiable on connected domains such as $\mathbb{R}^n$: overlapping $\epsilon$-balls form chains that connect any two inputs, and force a robust classifier to be constant. 
Existing work has addressed this limitation by weakening the guarantee: 
local robustness concerns only individual inputs~\cite{li2024certifying}; confidence-based notions restrict robustness to high-confidence
inputs~\cite{athavale2024verifying,leino2021globally}; probabilistic global robustness~\cite{mangal2019robustness,levy2023groma}.
However, probabilistic notions constrain only the density of counterexamples, not their geometric structure (see Fig.~\ref{fig:grid}).

To address this limitation, we introduce $(\mu,\epsilon)$-robustness:
Every input must lie within distance $\mu$ of a same-label point at which the classifier is locally $\epsilon$-robust. 
The definition recovers $\epsilon$-robustness at $\mu = 0$ and, unlike probabilistic
relaxations, imposes an explicit coverage \emph{structure} on non-robust points.
Two of our main results, a covering characterization of robustness and a necessary-and-sufficient criterion for the
existence of a non-constant robust classifier, hold over arbitrary semimetric spaces,
providing a domain-independent foundation.
Further, the impossibility result above is a consequence of \emph{connected
continuous} domains, whereas deployed classifiers compute over finite-precision
arithmetic. In fact, verification methods reasoning over reals can be unsound 
due to their abstraction~\cite{jia2021exploiting,cordeiro2025neural}. 
Since deployed systems operate over finite representable values rather than continuous domains, robustness becomes a quantitative question.
We show that the existence of a non-constant robust classifier is determined exactly 
by the relationship between $\epsilon$ ($\mu$) and the gaps between consecutive representable values. 

\paragraph{Contributions.}
Studying $(\mu,\epsilon)$-robustness gives us domain-independent structural results and quantitative guarantees for finite-precision arithmetic:
\begin{enumerate}
\item \textbf{Covering characterization.}
We introduce a covering characterization of $(\mu,\epsilon)$-robustness (Section~\ref{sec:partitions-covering}), extending the classical connection between $\epsilon$-robustness and separability. In the distribution-free setting, we establish separability as a necessary condition for global robustness~\cite{yang2020closer,mangal2023certifying}.

\item \textbf{Existence criteria.}
We derive necessary and sufficient conditions for the existence of a non-constant robust classifier (Section~\ref{sec:existence}), both for arbitrary semimetric spaces and for finite totally ordered domains.

\item \textbf{Precision-dependent bounds.}
We obtain bounds on $\epsilon$ and $\mu$ in terms of precision (ULP) metrics (Section~\ref{sec:upper-bounds}). These bounds identify the largest representable gap as the key quantity governing robustness, explain empirically observed robustness gains under coarser quantization~\cite{xiao2023robustmq,shrestha2024properties,duncan2020relative}, and extend to some partially ordered domains.
\end{enumerate}
Our analysis characterizes \emph{when} robust classifiers exist. 
Whether an \emph{accurate} robust classifier exists, and thus the robustness-accuracy
tradeoff~\cite{zhang2019theoretically}, is orthogonal to our focus. 
Along the way, our framework provides a unified, non-probabilistic perspective on several results 
concerning global robustness that have previously appeared in disparate parts of the literature.

%\begin{enumerate}
%    \item In Section~\ref{sec:new-robustness}, we define $(\mu,\epsilon)$-robustness and discuss alternative formulations of global classifier robustness.
%    \item In Section~\ref{sec:partitions-covering}, we formalize $(\mu,\epsilon)$-robustness on induced by classifiers partitions and show its equivalence to a covering property.
%    \item In Section~\ref{sec:existence}, we establish necessary and sufficient existence criteria for a non-constant robust classifier in general case and on finite totally ordered semimetric spaces, generalizing a line of finite precision numbers.
%    \item In Section~\ref{sec:upper-bounds}, we present upper bounds on $\epsilon$ and $\mu$ for non-constant classifiers on finite totally ordered semimetric spaces as well as their products.
%\end{enumerate}